\section{Experimental Setup}
\label{sec:setup}

\textbf{Datasets.}  \autoref{tab:datasets} summarizes the three
flights used for characterization and the main evaluation, recorded
as rosbags in a Vicon \ac{mocap} lab on the platform described in
\autoref{sec:overview}.  Additional held-out flights
(\autoref{sec:heldout}) and four public ICINS-2021 flights with
VINS pseudo ground truth (\autoref{sec:baselines}) test transfer
beyond these recordings.

\begin{table}[t]
\centering
\caption{Dataset Summary}
\label{tab:datasets}
\small
\begin{tabular}{@{}lcccc@{}}
\toprule
Flight & Dur.\ (s) & $v_{\max}$ (m/s) & $|\omega|_{\max}$ (rad/s) & Frames \\
\midrule
Slow racing & 25.8 & 3.65 & 2.62 & 258 \\
Fast racing & 18.0 & 5.66 & 9.14 & 179 \\
Backflips\textsuperscript{*} & 18.1 & $\sim$4.3 & $\sim$10.9 & 181 \\
\bottomrule
\multicolumn{5}{@{}p{\columnwidth}@{}}{\scriptsize \textsuperscript{*}Backflip
speed and angular rate are 95th percentiles rather than maxima because
\ac{mocap} tracking artifacts dominate the instantaneous peaks.}
\end{tabular}
\end{table}

\textbf{Configuration.}  Unless varied in an ablation, all flights
on both platforms use the configuration of \autoref{sec:method}:
a 3\,s window, 0.3\,s stride, and knot spacings
$\Delta t_p=40$\,ms and $\Delta t_o=16$\,ms, with no per-flight
retuning.  Radar extrinsics are fixed at each platform's
calibration, measured on ours and published with the public
dataset.  Nothing else was re-measured for the second platform: the noise-model
constants and stream weights are those measured on our own sensor and
flights (\autoref{sec:char}); the per-sample \ac{imu} weights were
applied unchanged at its 409\,Hz rate, and \autoref{sec:baselines}
reports what that costs.  Initialization uses radar and \ac{imu}
alone; heading aiding uses reference yaw as an idealized
magnetometer surrogate with $\lambda_\psi=0.6$ (\autoref{sec:yaw}).

\textbf{Evaluation.}  The reference is \ac{mocap} on our flights and the
public dataset's pseudo ground truth on theirs.
Estimates are aligned to the reference by the
rigid transform at the first frame only, without scale correction or
a whole-trajectory fit.  Such a fit is common in odometry work and is the
weaker protocol, since it can absorb accumulated drift into the transform.
The final 3\,s is excluded from evaluation.
Velocity \ac{rmse} compares the spline's analytical derivative with
the low-pass filtered reference velocity.  \emph{Position drift}
denotes position \ac{rmse} divided by the reference path length,
expressed as a percentage.  The \emph{live} \ac{rmse} uses each
window's newest stride before subsequent refinement; the
\emph{settled} \ac{rmse} uses the trajectory after all windows have
refined it.  For comparison with published results,
\autoref{tab:baselines}(b) additionally reports the whole-trajectory
\ac{ate} under Umeyama alignment.